\section{Rigidity of the PR-I Radial and Compact-Boundary Block}
\label{sec:pr35-pr1-radial-compact-rigidity}

The first uniqueness problem is upstream of every electromagnetic,
electroweak, flavor, and radiative result.  It asks whether the PR-I radial
operator and compact boundary map are fixed by the declared projection rules,
or whether their published values select one member of an unreported family.
This section proves rigidity in an explicit admissible class and separates the
exact result from the remaining geometric assumptions.

\subsection{Minimal Radial Candidate Class}
\label{subsec:pr35-minimal-radial-class}

Let
\begin{equation}
\mathfrak V_4
=
\left\{
V(w)=\Lambda_X^2\left(1+a_2w^2+a_4w^4\right)
\;\middle|\;
a_2>0,\ a_4>0
\right\}
\label{eq:pr35-radial-admissible-class}
\end{equation}
be the minimal even quartic class.  Unit local radial stiffness gives
\begin{equation}
\left.
\frac{1}{2\Lambda_X^2}\frac{d^2V}{dw^2}
\right|_{w=0}
=a_2=1,
\label{eq:pr35-a2-unique}
\end{equation}
whereas the spatial projection trace gives
\begin{equation}
a_4
=
\frac{\operatorname{Tr}P_{\rm space}}
{\operatorname{Tr}I_{3+1}}
=
\frac34.
\label{eq:pr35-a4-unique}
\end{equation}
The two coefficient constraints are independent; their coefficient matrix has
determinant $4\neq0$.  Consequently, there is exactly one member of
\(\mathfrak V_4\) satisfying both constraints:
\begin{equation}
\boxed{
V_{X,\star}(w)
=
\Lambda_X^2\left(1+w^2+\frac34w^4\right).}
\label{eq:pr35-radial-unique}
\end{equation}
Furthermore,
\begin{equation}
\frac{V_{X,\star}''(w)}{\Lambda_X^2}
=2+9w^2>0,
\qquad
\lim_{|w|\to\infty}V_{X,\star}(w)=+\infty.
\end{equation}
Thus the associated one-dimensional Schr\"odinger operator is confining and has a lower-bounded discrete simple spectrum on its standard self-adjoint
domain \cite{ReedSimon_1975}.

Equation~\eqref{eq:pr35-radial-unique} is unique in the class \(\mathfrak V_4\).  It is not an unrestricted functional-uniqueness theorem. If higher even or nonpolynomial terms are admitted without additional constraints, an infinite family remains.  The minimal-even-quartic condition is therefore part of the theorem domain.

\subsection{Interior Recurrence as a Numerical Semigroup}
\label{subsec:pr35-compact-semigroup}

The primitive compact returns have lengths $3$ and $4$.  Every recurrent
interior length belongs to
\begin{equation}
S_{3,4}
=
\langle3,4\rangle
=
\{3a+4b\mid a,b\in\mathbb Z_{\ge0}\}.
\label{eq:pr35-semigroup}
\end{equation}
The positive gaps of this semigroup are exactly
\begin{equation}
\mathbb Z_{>0}\setminus S_{3,4}
=
\{1,2,5\}.
\label{eq:pr35-semigroup-gaps}
\end{equation}
Indeed, $6=3+3$, $7=3+4$, and $8=4+4$ cover all residue classes modulo $3$; adding copies of $3$ proves that every integer greater than or equal to $6$ lies in $S_{3,4}$. A finite-rank boundary word must contain the leave--return pair $A^2$ followed by the primitive spatial return $A^3$.  It therefore has exponent at least $5$.  Equation~\eqref{eq:pr35-semigroup-gaps} then proves that
\begin{equation}
XA^5=XA^{2+3}
\label{eq:pr35-unique-positive-terminal}
\end{equation}
is the unique positive boundary word that closes and is not already generated by the recurrent interior monoid. The first repeated full-trace return is $A^{4+4}=A^8$.  Under the PR finite-rank orientation rule---retain the spatial primitive channel and
subtract the first double full-trace recirculation outside the finite slice---the signed exclusion is $-XA^8$.  This exponent and sign are unique under that orientation rule.  The semigroup alone does not derive the rule; the rule remains an explicit geometric condition of the present theorem.

The resulting signed terminal set is
\begin{equation}
\boxed{
\operatorname{Adm}_{\Pi_{13}}
=
\{\epsilon,A^3,A^4,XA^5,-XA^8\}.}
\label{eq:pr35-signed-terminal-set}
\end{equation}

\subsection{Exact Cofactor Closure}
\label{subsec:pr35-exact-cofactor}

With $T_A=1/4$ and $T_X=3/4$,
\begin{equation}
M_A
=
\begin{pmatrix}
T_A^3&T_A^4\\
1&0
\end{pmatrix}
\end{equation}
gives
\begin{equation}
D_A
=
\det(I-M_A)
=
1-T_A^3-T_A^4
=
\frac{251}{256}.
\label{eq:pr35-da-exact}
\end{equation}
The signed terminal set gives
\begin{equation}
N_{\rm bc}
=
1+T_A^3+T_A^4+T_X(T_A^5-T_A^8)
=
\frac{267453}{262144}.
\label{eq:pr35-nbc-exact}
\end{equation}
Consequently,
\begin{equation}
c_{\rm bc}
=
T_X\left[1+T_A^2-T_A^6\frac{N_{\rm bc}}{D_A}\right]
=
\frac{3354902985}{4211081216}
=
0.7966844648478990532\ldots .
\label{eq:pr35-cbc-exact}
\end{equation}
These are rational identities and contain no empirical input.

\subsection{Minimal Finite-rank Projector}
\label{subsec:pr35-pi13-uniqueness}

Let the odd radial eigenbasis be \(\{\phi_1,\phi_3,\phi_5,\ldots\}\).  Require a coordinate spectral projector that contains both the fundamental odd winding-linked channel \(\phi_1\) and the spatial primitive-return channel \(\phi_3\), preserves
parity, and has minimal rank.  There is exactly one rank-two coordinate projector satisfying these requirements:
\begin{equation}
\boxed{
\Pi_{13}
=
|\phi_1\rangle\langle\phi_1|
+
|\phi_3\rangle\langle\phi_3|.}
\label{eq:pr35-pi13-unique}
\end{equation}
This theorem is conditional on the identification of channels (1) and (3) as required.  A derivation excluding every additional odd channel is a stronger statement than minimal-projector uniqueness.

\subsection{Two-mode Spectral Isolation and Leakage}
\label{subsec:pr35-two-mode-leakage}

Define
\begin{equation}
K_X
=
\frac{O_X-\lambda_0}{\lambda_1-\lambda_0}
\end{equation}
and
\begin{equation}
P_{\rm geom}^{\rm bc}
=
\Pi_{13}K_X^2
\left(c_{\rm bc}+w^2+\frac34w^4\right)
K_X^2\Pi_{13}.
\label{eq:pr35-pgeom}
\end{equation}
On the range of \(\Pi_{13}\), Eq.~\eqref{eq:pr35-pgeom} is a real symmetric $2\times2$ matrix.  Its reconstructed off-diagonal element is nonzero and its eigenvalues are distinct.  The dominant normalized eigenvector is therefore unique up to an overall sign.  The sign cancels from
\begin{equation}
p_1
=
|\langle\phi_1|\Omega_{\rm bc}\rangle|^2,
\qquad
p_3=1-p_1,
\end{equation}
so the leakage probabilities are unique in the two-mode block. The independent harmonic-oscillator-basis certificate gives
\begin{align}
\lambda_0&=2.322872581463\ldots,\\
\lambda_1&=5.375889248196\ldots,\\
\lambda_3&=13.233893408730\ldots,\\
\mu_{\min}^2&=3.053016666733\ldots,\\
p_1&=7.68534695089\times10^{-4}.
\end{align}
No empirical electromagnetic constant enters this calculation.  The compact
stiffness and tree normalization are generated as
\begin{align}
q_{\rm bc}
&=
p_1(\lambda_1-\lambda_0)
+(1-p_1)(\lambda_3-\lambda_0)
=10.904981678435\ldots,\\
\alpha_{\rm PR,tree}^{-1}
&=4\pi q_{\rm bc}
=137.036041314016\ldots .
\label{eq:pr35-alpha-tree-regenerated}
\end{align}

This calculation reproduces the frozen PR-III high-resolution branch.  It
does not reproduce the earlier finite-resolution PR-I value
\(p_1=7.6528903366\times10^{-4}\).  The latter is therefore classified as
superseded finite-resolution data rather than as the canonical operator limit.

\subsection{Claim Boundary}
\label{subsec:pr35-pr1-rigidity-boundary}

The radial coefficient solution, numerical-semigroup gap theorem, positive
terminal word, rational cofactor arithmetic, and minimal coordinate projector
are exact in the declared class.  The two-mode leakage and high-resolution
compact stiffness are numerically isolated and reproducible.  The remaining
conditional inputs are the minimal-quartic function class, the primitive return
set \(\{A^3,A^4\}\), the required-channel identification
\(\{\phi_1,\phi_3\}\), and the finite-rank orientation rule producing
\(-XA^8\).  These conditions define the precise boundary between the proved
PR-I rigidity result and the next geometric uniqueness obligation.
